\documentclass[../paper.tex]{subfiles}
\begin{document}
\section{Beyond $2\times2$}\label{sec:beyond2x2}
In this section I briefly explore the dynamics of commitment beyond $2\times 2$ games. Consider the game in table \ref{tab:3x3}. This game has the unique Nash equilibrium $\{A_1,A_1\}$. Furthermore note that for player 1, $A_2$ is strictly dominated by $A_1$, and for player 2, $A_3$ is strictly dominated by $A_1$.
\begin{table}[h!]
\centering
\begin{tabular}{c|ccc}
    & $A_1$ & $A_2$ & $A_3$ \\
    \hline
    $A_1$ & $1,1$ & $3,0$ & $\epsilon,0$ \\
    $A_2$ & $0,\epsilon$ & $2,1$ & $0,0$ \\
    $A_3$ & $0,3$ & $0,0$ & $1,2$ \\
\end{tabular}
\caption{A $3\times3$ game }
\label{tab:3x3}
\end{table}
When commitment is introduced, there are two pure strategy subgame perfect Nash equilibria: 1 commits to $A_2$ and 2 best responds with $A_2$, or 2 commits to $A_3$ and 1 best responds with $A_3$. Players now choose to commit to strategies that are strictly dominated in the base game \parencite[][also make the observation that commitment to dominated strategies is possible]{harrenstein_commitment_2007}. 
Furthermore, there is an interior mixed equilibrium, where player 1 commits to $A_2$ with probability $\frac{1}{2}$ and waits with probability $\frac{1}{2}$ and player 2 commits to $A_3$ with probability $\frac{1}{2}$ and waits with probability $\frac{1}{2}$. Again, players make probabilistic commitments to strictly dominated strategies. The resulting expected payoff of this interior mixed equilibrium is $1$. 
 
Next, modify only the payoffs that the \textit{opponent} receives when $A_2$ or $A_3$ is played, as in Table \ref{tab:3x3modified}. Choosing $A_2$ still gives player 1 the same payoffs, but now imposes a large cost on player 2, and vice versa for $A_3$. The unique base-game equilibrium $\{A_1,A_1\}$ is unchanged, still giving each player a payoff of 1. With commitment, the game's mixed equilibrium becomes interesting: players still mix equally between committing $A_2$ and $W$, or committing $A_3$ and $W$, respectively. The expected payoff in this case drops drastically to $1 - \frac{10^9}{2}$ for each player, another demonstration of how commitment can be mutually harmful.
\begin{table}[h!]
\centering
\begin{tabular}{c|ccc}
    & $A_1$ & $A_2$ & $A_3$ \\
    \hline
    $A_1$ & $1,1$ & $3,0$ & $\epsilon -\textcolor{red}{10^{9}},0$ \\
    $A_2$ & $0,\epsilon-\textcolor{red}{10^{9}}$ & $2,1-\textcolor{red}{10^{9}}$ & $0-\textcolor{red}{10^{9}},0-\textcolor{red}{10^{9}}$ \\
    $A_3$ & $0,3$ & $0,0$ & $1-\textcolor{red}{10^{9}},2$ \\
\end{tabular}
\caption{A modified $3\times3$ game}
\label{tab:3x3modified}
\end{table}
\newline
The games above illustrate that, in contrast to $2\times2$ games, $3\times3$ games can have mixed equilibria with probabilistic commitment even if the base game only has one Nash equilibrium. 
Furthermore, while in the $2 \times 2$ case payoffs could decrease relative to the base game equilibrium, the difference in payoffs could not become too large, since both were functions of the same payoff parameters. By contrast, in the $3 \times 3$ case, this difference is less restricted and can grow arbitrarily large, as demonstrated.
Overall, this section illustrated just one possible dynamic beyond $2\times2$ games. A more systematic exploration of such games would likely reveal further interesting phenomena, but lies beyond the scope of this thesis.

\end{document}